\chapter{Sesión 4: Extracción de características, Bolsa de Palabras Visuales, Detección de Movimiento y Seguimiento de Objetos}
\label{chapter:introduction_ses_4}

\section{Materiales}

En esta práctica se trabajará con los siguientes recursos (puede encontrarlos en la sección de Moodle \textit{Laboratorio/Sesión 4}):

\begin{itemize}
    \item \textbf{\texttt{partA.ipynb}}: notebook con el código que deberá desarrollar.
    \item \textbf{\texttt{partB.ipynb}}: notebook con el código que deberá desarrollar.
    \item \textbf{\texttt{partC.ipynb}}: notebook con el código que deberá desarrollar.
    \item \textbf{data}: carpeta con el material para trabajar durante la práctica.
\end{itemize}

\section{Apartados de la práctica}

La Sesión 4 del laboratorio está dividida en los siguientes apartados:

\begin{itemize}
    \item Apartado A: Detección de líneas rectas.
    \item Apartado B: Detección de puntos de interés y bolsa de palabras.
    \item Apartado C: Sustracción de Fondo.
    \item Apartado D: Flujo Óptico.
    \item Apartado E: Filtro de Kalman para Seguimiento de Objetos.

\end{itemize}

\section{Observaciones}

Aunque el guion de la práctica y los comentarios en Markdown del Notebook estén escritos en español, observe que todo aquello que aparece en las celdas de código está escrito en inglés. Es una buena práctica que todo su código esté escrito en inglés.

Aquellas partes del código que deberá completar están marcadas con la etiqueta \textbf{\texttt{TODO}}.

Es muy importante que trabaje consultando la documentación de OpenCV\footnote{\href{https://docs.opencv.org/4.x/index.html}{Documentación de OpenCV}: \url{https://docs.opencv.org/4.x/index.html}} para familiarizarse de cara al examen. Tenga en cuenta que en los exámenes no podrá utilizar herramientas de ayuda como Copilot.

\section{Qué va a aprender}

Al finalizar esta práctica, aprenderá a implementar algoritmos de detección de líneas rectas, puntos de interés y a usar de bolsa de palabras visuales. Así como a implementar algoritmos detección de movimiento mediante sustracción de fondo y flujo óptico, y explorará el filtro de Kalman para el seguimiento de objetos.

\section{Evaluación}

La nota que obtenga en esta sesión de laboratorio será la misma que obtenga su pareja. Los apartados de la práctica serán evaluados como refleja la Tabla \ref{table:evaluacion}.

\begin{table}[h!]
    \centering
    \begin{tabular}{|c|c|c|}
    \hline
    \textbf{Tarea} & \textbf{Valor} & \textbf{Resultado} \\
    \hline
    Pregunta A.1 & 1.0 & \\
    \hline
    Pregunta B.1 & 1.5 & \\
    \hline
    Pregunta B.2 & 1.5 & \\
    \hline
    Pregunta C.1 & 1.0 & \\
    \hline
    Pregunta C.2 & 1.0 & \\
    \hline
    Pregunta D.1 & 1.0 & \\
    \hline
    Pregunta D.2 & 1.0 & \\
    \hline
    Pregunta E.1 & 1.0 & \\
    \hline
    Pregunta E.2 & 1.0 & \\
    \hline
    \textbf{Total} & \textbf{10.0} & \\
    \hline
    \end{tabular}
    \caption{Valoración de los apartados de la práctica.}
    \label{table:evaluacion}
\end{table}
